\section{The \texttt{Filter\_gen} McStas Component}
Release: McStas 1.6

This components may either set the flux or change it (filter-like), using
an external data filename.

\subsection*{Identification}
\begin{itemize}
  \item \textbf{Author:} E. Farhi
  \item \textbf{Origin:} ILL
  \item \textbf{Date:} Dec, 15th, 2002
\end{itemize}

\subsection*{Description}
This component changes the neutron flux (weight) in order to match a reference table in a filename. Typically you may set the neutron flux (source-like), or multiply it using a transmission table (filter-like). The component may be placed after a source, in order to e.g. simulate a real source from a reference table, or used as a filter (BeO) or as a window (Al). The behaviour of the component is specified using the 'options' parameter, or from the filename itself (see below) If the thickness for the transmission data filename D was t0, and a different thickness t1 would be required, then the resulting transmission is:

\begin{verbatim}
D^(t1/t0).
\end{verbatim}

You may use the 'thickness' and 'scaling' parameter for that purpose.

\textbf{File format:} This filename may be of any 2-columns free format (k[\AA{}-1],p), (omega[meV],p) and (lambda[\AA{}],p) where p is the weight. The type of the filename may be written explicitely in the filename, as a comment, or using the 'options' parameter. Non mumerical content in filename is treated as comment (e.g. lines starting with '\#' character). A table rebinning and linear interpolation are performed.

EXAMPLE : in order to simulate a PG filter, using the lib/data/HOPG.trm file Filter\_gen(xwidth=.1 yheight=.1, filename="HOPG.trm") A Sapphire filter, using the lib/data/Al2O3\_sapphire.trm file Filter\_gen(xwidth=.1 yheight=.1, filename="Al2O3\_sapphire.trm") A Berylium filter, using the lib/data/Be.trm file Filter\_gen(xwidth=.1 yheight=.1, filename="Be.trm") an other possibility to simulate a Be filter is to use the PowderN component: PowderN(xwidth=.1, yheight=.1, zdepth=.1, reflections="Be.laz", p\_inc=1e-4)

in this filename, the comment line \# wavevector multiply sets the behaviour of the component. One may as well have used options="wavevector multiply" in the component instance parameters.

\subsection*{Input parameters}
Parameters in \textbf{boldface} are required; the others are optional.

\begin{longtable}{p{0.22\textwidth}p{0.12\textwidth}p{0.46\textwidth}p{0.14\textwidth}}
\toprule
\textbf{Name} & \textbf{Unit} & \textbf{Description} & \textbf{Default} \\
\midrule
\endhead
filename & str & name of the filename to look at (first two columns data). Data D should rather be sorted (ascending order) and monotonic filename may contain options (see below) as comment & 0 \\
options & str & string that can contain: "[ k p ]"      or "wavevector" for filename type, "[ omega p]"   or "energy", "[ lambda p ]" or "wavelength", "set"          to set the weight according to the table,"multiply"     to multiply (instead of set) the weight by factor,"add"          to add to current flux,"verbose"      to display additional informations. & 0 \\
xmin & m & dimension of filter & -0.05 \\
xmax & m & dimension of filter & 0.05 \\
ymin & m & dimension of filter & -0.05 \\
ymax & m & dimension of filter & 0.05 \\
xwidth & m & Width/diameter of filter). Overrides xmin,xmax. & 0 \\
yheight & m & Height of filter. Overrides ymin,ymax. & 0 \\
thickness & 1 & relative thickness. D = D\textasciicircum{}(thickness). & 1 \\
scaling & 1 & scaling factor. D = D*scaling. & 1 \\
verbose & 1 & Flag to select verbose output. & 0 \\
\bottomrule
\end{longtable}

\subsection*{Links}
\begin{itemize}
  \item Component source code found in file \texttt{Filter\_gen.comp}.
  \item \htmladdnormallink{HOPG.trm}{../data/HOPG.trm} filename as an example.
\end{itemize}
\IfFileExists{optics/Filter_gen_static.tex}{\subsection*{Example transmission data files}

Some example transmission data files are provided with \MCS in the
\verb+MCSTAS/data+ directory as \verb+*.trm+ files, for use with the
\texttt{filename} parameter above:

\begin{table}[htb!]
  \begin{center}
  {\let\my=\\
    \begin{tabular}{|l|p{0.7\textwidth}|}
    \hline
    File name & Description \\
    \hline
    Be.trm       & Berylium filter for cold neutron spectrometers (e.g. $k < 1.55$  \AA$^{-1}$) \\
    HOPG.trm     & Highly oriented pyrolithic graphite for $\lambda/2$ filtering (e.g. thermal beam at $k$ = 1.64, 2.662, and 4.1 \AA$^{-1}$) \\
    Sapphire.trm & Sapphire (Al$_2$O$_3$) filter for fast neutrons ($k > 6$ \AA$^{-1}$) \\
    \hline
    \end{tabular}
    \caption{Some transmission data files to be used with the Filter\_gen component.}
    \label{t:filter-transmission-data}
  }
  \end{center}
\end{table}
}{}